\chapter{Kernel Jupyter}

\section{O que é o kernel Jupyter}

O \hayashi{} inclui um kernel Jupyter nativo, chamado \texttt{hay-kernel}. Ele
permite executar código \hayashi{} dentro de notebooks Jupyter, do JupyterLab ou
de qualquer frontend compatível com o protocolo Jupyter.

Diferente do REPL, o kernel mantém o estado entre células: variáveis, funções e
dados definidos em uma célula ficam disponíveis nas células seguintes do mesmo
notebook.

\section{Compilação e instalação}

O kernel faz parte do repositório \hayashi{} e é compilado com a feature
\texttt{native}:

\begin{lstlisting}[style=inline, language=sh]
cd hayashi
cargo build --release --bin hay-kernel --features native
\end{lstlisting}

Depois, instale o especificação do kernel:

\begin{lstlisting}[style=inline, language=sh]
./target/release/hay-kernel --install
\end{lstlisting}

Isso registra o kernel em \texttt{\$HOME/.local/share/jupyter/kernels/hayashi/}
(ou o diretório equivalente na plataforma).

Verifique se o kernel aparece na lista:

\begin{lstlisting}[style=inline, language=sh]
jupyter kernelspec list
\end{lstlisting}

A saída deve conter uma linha como:

\begin{lstlisting}[style=inline]
hayashi    ~/.local/share/jupyter/kernels/hayashi
\end{lstlisting}

\section{Usando o kernel}

\subsection{JupyterLab}

Inicie o JupyterLab:

\begin{lstlisting}[style=inline, language=sh]
jupyter lab
\end{lstlisting}

Crie um novo notebook e selecione o kernel \textbf{Hayashi} no launcher ou no menu
de kernel.

\subsection{Jupyter Notebook}

Inicie o Jupyter Notebook:

\begin{lstlisting}[style=inline, language=sh]
jupyter notebook
\end{lstlisting}

Crie um novo notebook e escolha \textbf{Hayashi} na lista de kernels.

\subsection{Executando células}

As células aceitam código \hayashi{} diretamente:

\begin{lstlisting}[caption={Exemplo de célula no Jupyter}]
let x = 1 + 2
print(x)
\end{lstlisting}

A saída \texttt{3} é exibida logo abaixo da célula. Variáveis persistem entre
células:

\begin{lstlisting}[caption={Célula posterior com a mesma sessão}]
let y = x * 10
print(y)
\end{lstlisting}

\subsection{Jupyter Console}

Se estiver instalado, o console também pode ser usado:

\begin{lstlisting}[style=inline, language=sh]
jupyter console --kernel=hayashi
\end{lstlisting}

\section{Desinstalação}

Remova o diretório do especificação do kernel:

\begin{lstlisting}[style=inline, language=sh]
rm -rf ~/.local/share/jupyter/kernels/hayashi
\end{lstlisting}

\section{Limitações em ambientes hospedados}

Google Colab, Kaggle Notebooks e ambientes similares não permitem registrar
kernels personalizados de forma simples e o runtime é efêmero. Nessas situações,
a alternativa prática é subir um binário pré-compilado \texttt{hay} e chamá-lo de
células Python, eventualmente com uma magic cell como \texttt{\%\%hay}.

Veja o capítulo \ref{chap:colab} para um exemplo de gambiarra funcional.
